\documentclass[11pt,a4paper]{article}
\usepackage[T1]{fontenc}
\usepackage[utf8]{inputenc}
\usepackage[english]{babel}
\usepackage{amsmath,amssymb}
\usepackage[margin=2cm]{geometry}
\setlength{\parindent}{0pt}
\setlength{\emergencystretch}{3em}
\allowdisplaybreaks

\title{Parity, remainders and divisibility (14)}
\author{}
\date{}
\begin{document}
\maketitle
$\mathbb{N} = \{1, 2, 3, \ldots\}$, $\mathbb{N}_{0} = \mathbb{N} \cup \{0\}$. Every statement below is true: each says that an expression in an integer, divided by a fixed number, leaves one of the listed remainders.

\begin{enumerate}
\item $\displaystyle \forall x \in \mathbb{Z}\ \exists y \in \mathbb{Z}\ (x = 2 y \lor x = 2 y + 1)$
\item $\displaystyle \forall x \in \mathbb{Z}\ \exists y \in \mathbb{Z}\ (x = 3 y \lor x = 3 y + 1 \lor x = 3 y + 2)$
\item $\displaystyle \forall x \in \mathbb{Z}\ \exists y \in \mathbb{Z}\ (x^{2} + x = 2 y)$
\item $\displaystyle \forall x \in \mathbb{Z}\ \exists y \in \mathbb{Z}\ (x^{3} - x = 6 y)$
\item $\displaystyle \forall x \in \mathbb{Z}\ \exists y \in \mathbb{Z}\ (x^{2} = 4 y \lor x^{2} = 4 y + 1)$
\item $\displaystyle \forall x \in \mathbb{Z}\ \exists y \in \mathbb{Z}\ (x^{2} = 3 y \lor x^{2} = 3 y + 1)$
\item $\displaystyle \forall x \in \mathbb{N}\ \exists y \in \mathbb{N}_{0}\ (x = 2 y + 1 \lor x = 2 y + 2)$
\item $\displaystyle \forall x \in \mathbb{Z}\ \exists y \in \mathbb{Z}\ (x^{3} = 9 y \lor x^{3} = 9 y + 1 \lor x^{3} = 9 y - 1)$
\item $\displaystyle \forall x \in \mathbb{Z}\ \exists y \in \mathbb{Z}\ (x^{5} - x = 5 y)$
\item $\displaystyle \forall x \in \mathbb{Z}\ \exists y \in \mathbb{Z}\ (x \left(x + 1\right) \left(2 x + 1\right) = 6 y)$
\item $\displaystyle \forall x \in \mathbb{N}_{0}\ \exists y \in \mathbb{N}_{0}\ (x^{2} + 3 x = 2 y)$
\item $\displaystyle \forall x \in \mathbb{Z}\ \exists y \in \mathbb{Z}\ (x^{2} + x + 1 = 2 y + 1)$
\item $\displaystyle \forall a \in \mathbb{Z}\ \exists b \in \mathbb{Z}\ (a^{4} = 5 b \lor a^{4} = 5 b + 1)$
\item $\displaystyle \forall x \in \mathbb{N} \setminus \{1\}\ \exists y \in \mathbb{N}\ (x = 2 y \lor x = 2 y + 1)$
\end{enumerate}
\end{document}
